\chapter{exec、ELF 装载与用户栈}

\section{程序、进程与程序映像}

\code{fork} 复制了一个进程的执行环境，而 \code{exec} 把当前进程的用户态程序替换成另一个程序。这个替换不是“创建新进程”，因为 pid、父子关系、部分资源限制和很多内核对象仍然属于同一个进程；它也不是普通文件读取，因为内核必须验证可执行文件格式、建立新的地址空间、设置入口点、构造用户栈、处理解释器、重置部分信号语义并关闭带 \code{close-on-exec} 标志的 fd。

ELF 装载器是系统调用边界、文件系统、虚拟内存和权限模型的交汇点。用户给内核一个路径和参数数组，内核必须在不信任用户内存的前提下读取这些字符串；路径解析得到 inode 后，还要检查执行权限、文件类型、挂载选项和安全策略；读取 ELF 头后，装载器根据 program header 而不是 section header 建立映射。section 是链接和调试视角，program segment 才是运行时装载视角。

\begin{keyidea}
\code{exec} 的核心契约是原子替换用户态映像：成功后旧程序不能继续执行；失败时旧程序必须仍然可运行，除非失败发生在明确不可恢复的提交点之后。
\end{keyidea}

真实系统还要处理动态链接。若 ELF 带 \code{PT\_INTERP}，内核并不直接跳到主程序入口，而是同时映射解释器，把控制权交给动态链接器。动态链接器再加载共享库、重定位符号，最后跳到用户程序的入口。内核只负责把 aux vector、入口、program header 信息和随机数等必要元数据放到用户栈。

\section{exec 引起的完整状态替换}

一个教学内核中的 \code{execve(path, argv, envp)} 可以拆成九个阶段：

\begin{enumerate}
\tightlist
\item 从用户内存复制 path、argv、envp，设置长度上限。
\item 路径解析并打开可执行 inode。
\item 检查文件类型、执行权限、noexec 挂载选项和安全策略。
\item 读取并验证 ELF header。
\item 扫描 program header，计算新地址空间布局。
\item 创建临时地址空间，映射 \code{PT\_LOAD} 段和用户栈。
\item 若存在 \code{PT\_INTERP}，装载解释器并设置实际入口。
\item 构造用户栈：argc、argv 指针、envp 指针、auxv、字符串和对齐填充。
\item 提交：切换进程 mm、寄存器入口和栈指针，关闭 \code{CLOEXEC} fd。
\end{enumerate}

这里最容易讲错的是失败回滚。前七步应尽量在临时对象中完成：临时 mm、临时 VMA、临时页、临时 fd 引用。只有全部验证成功后，才进入提交阶段并销毁旧地址空间。否则一个坏 ELF 就可能把调用者进程破坏到无法返回错误码。

\begin{longtable}{p{0.18\linewidth}p{0.36\linewidth}p{0.34\linewidth}}
\toprule
对象 & 成功后状态 & 失败回滚 \\
\midrule
旧 mm & 被新 mm 替换并释放引用 & 保持可运行 \\
新 mm & 成为 current 的地址空间 & 释放所有 VMA 和页 \\
文件引用 & 可执行文件可关闭或保留映射引用 & put inode/file 引用 \\
fd table & 关闭 \code{FD\_CLOEXEC} 条目 & 不应改变 \\
信号处理 & 捕获处理器重置，忽略语义按规则保留 & 不应改变 \\
寄存器 & RIP 指向入口，RSP 指向新用户栈 & 不应改变 \\
\bottomrule
\end{longtable}

把这条路径放进一个具体命令会更清楚。假设 shell 进程 pid 为 3182，执行：

\begin{lstlisting}
execve("/usr/bin/grep",
       ["grep", "-n", "main", "a.txt"],
       envp);
\end{lstlisting}

进入内核时，\code{path}、\code{argv} 和 \code{envp} 都还在旧地址空间里。这里的旧地址空间不是“快要没用的垃圾”，而是内核目前唯一能读到这些字符串的地方。内核不能先释放旧 mm 再慢慢读参数；一旦旧页表被拆掉，\code{argv[2]} 指向的 \code{"main"} 可能就再也翻译不出来。因此第一轮工作是复制输入：先复制路径字符串，再按指针数组逐个复制 argv/envp 字符串，并统计总字节数。若 \code{argv[3]} 恰好跨到一个不可读页，正确结果是 \code{-EFAULT}，旧 shell 继续运行，fd table、signal handler 和旧地址空间都不能被改坏。

\begin{longtable}{p{0.16\linewidth}p{0.36\linewidth}p{0.36\linewidth}}
\toprule
阶段 & 内核手里的稳定证据 & 此时失败应怎样收场 \\
\midrule
复制用户参数 & 内核缓冲区里的 path、argv、envp 副本 & 返回 \code{EFAULT/E2BIG}，旧程序继续 \\
路径解析 & 指向 \code{/usr/bin/grep} 的 inode/file 引用 & 返回 \code{ENOENT/EACCES}，释放引用 \\
ELF 验证 & ELF header 和 program header 副本 & 返回 \code{ENOEXEC}，旧 mm 不变 \\
临时装载 & 新 mm、VMA、栈页、解释器信息 & 释放临时 mm，旧程序继续 \\
提交 & 新入口、新栈、新凭据、CLOEXEC 关闭集合 & 成功后不再回到旧用户代码 \\
\bottomrule
\end{longtable}

接着内核解析路径，得到稳定的 inode 和 file 引用。权限检查必须绑定这个对象，而不是只检查字符串 \code{"/usr/bin/grep"}。检查通过后，内核读取 ELF header。若这是动态链接的程序，program header 里通常有 \code{PT\_INTERP}，例如指向 \code{/lib64/ld-linux-x86-64.so.2}。这意味着“用户写的是执行 grep”，但第一条用户态指令很可能属于动态链接器；动态链接器根据 auxv 找到主程序的 program header，完成共享库装载和重定位，最后再跳到 grep 的入口。这个过程解释了为什么 exec 栈上要有 \code{AT\_PHDR}、\code{AT\_PHNUM}、\code{AT\_ENTRY}、\code{AT\_BASE} 和 \code{AT\_RANDOM}，它们不是装饰字段，而是用户态运行时启动所需的证据。

再看文件描述符。shell 可能提前打开了管道：

\begin{lstlisting}
grep -n main a.txt | less
\end{lstlisting}

这时 grep 进程继承的 \code{STDOUT\_FILENO} 可能不是终端，而是管道写端。exec 成功后，普通 fd 继续指向同一个 open file description；带 \code{FD\_CLOEXEC} 的 fd 则在提交时关闭。这个规则必须在提交点附近执行：太早关闭会导致 exec 失败后旧程序少了 fd；太晚关闭会让新程序短暂看到本不该继承的能力。于是 \code{FD\_CLOEXEC} 不是“关闭一些数字”这么简单，它是在程序映像替换边界上裁剪继承能力。

最后把新用户栈具体展开。假设栈顶从高地址向低地址增长，内核通常先把字符串放到栈下方，再压入指针表和 auxv。进入动态链接器或程序入口时，用户态看到的是：

\begin{lstlisting}
low address
  argc = 4
  argv[0] -> "grep"
  argv[1] -> "-n"
  argv[2] -> "main"
  argv[3] -> "a.txt"
  argv[4] = NULL
  envp[0] -> ...
  ...
  NULL
  auxv entries: AT_PHDR, AT_PHNUM, AT_ENTRY, AT_RANDOM, ...
  padding for alignment
  strings: "grep\0-n\0main\0a.txt\0..."
high address
\end{lstlisting}

如果这里指针顺序错了，\code{main(int argc, char** argv)} 看到的参数就会错；如果 16 字节对齐错了，某些 ABI 约定和运行时启动代码会在很早的位置崩溃；如果 \code{AT\_RANDOM} 指向了旧地址空间或未映射地址，libc 初始化 canary 时会 fault。用户栈构造不是格式化输出，而是内核和用户态启动代码之间的二进制契约。

\section{ELF 校验、段映射与用户栈}

\subsection{ELF header 验证}

ELF 验证要防止把任意文件当代码执行。最小检查包括 magic、class、endianness、machine、type、program header 范围和整数溢出。

\begin{lstlisting}
validate_elf(file):
  eh = read_at(file, 0, sizeof(Elf64_Ehdr))
  require eh.magic == 0x7f 'E' 'L' 'F'
  require eh.class == ELFCLASS64
  require eh.machine == EM_X86_64
  require eh.phentsize == sizeof(Elf64_Phdr)
  require eh.phnum <= MAX_PHDRS
  require checked_mul(eh.phnum, eh.phentsize, &ph_size)
  require checked_add(eh.phoff, ph_size, &ph_end)
  require ph_end <= file.size or file supports sparse read policy
\end{lstlisting}

复杂度是 \code{O(phnum)}。真正关键的是边界检查：\code{p\_offset + p\_filesz}、\code{p\_vaddr + p\_memsz}、页对齐和权限组合都可能溢出。装载器不能相信编译器或链接器一定生成合法文件，因为攻击者可以手工构造 ELF。

\subsection{段映射}

\code{PT\_LOAD} 段描述文件区间到虚拟地址区间的映射。文件大小 \code{p\_filesz} 是需要从文件读出的内容，内存大小 \code{p\_memsz} 包括 BSS 零填充。

\begin{lstlisting}
map_load_segment(mm, file, ph):
  require ph.memsz >= ph.filesz
  require same_page_offset(ph.vaddr, ph.offset)
  start = page_down(ph.vaddr)
  end   = page_up(ph.vaddr + ph.memsz)
  file_start = page_down(ph.offset)
  prot = prot_from_flags(ph.flags)
  map_file(mm, start, end, file, file_start, prot)
  zero_range(mm, ph.vaddr + ph.filesz, ph.vaddr + ph.memsz)
\end{lstlisting}

若采用 demand paging，\code{map\_file} 不必立即读入所有页面，只需建立 VMA，等 page fault 时按文件偏移填充页面。这样 \code{exec} 延迟低，内存占用也低。若采用 eager loading，实现简单但启动大程序会有明显 I/O 峰值。

用 \code{/usr/bin/grep} 的第一条指令看 demand paging。exec 提交完成后，寄存器 RIP 已经指向解释器或主程序入口，RSP 指向新用户栈；但入口所在的代码页可能还没有真正读入物理内存。CPU 返回用户态后开始取指，MMU 查入口虚拟地址的 PTE，发现 present=0，于是触发 instruction page fault。内核重新进入缺页处理，先用 fault 地址找到对应 VMA，再发现它来自可执行文件的 \code{PT\_LOAD} 段，权限是 read+exec，文件偏移由 \code{vaddr - segment.vaddr + segment.offset} 算出。若 page cache 里没有这一页，内核提交文件读；读完后安装只读可执行 PTE，返回用户态重试同一条取指。

\begin{longtable}{p{0.18\linewidth}p{0.36\linewidth}p{0.34\linewidth}}
\toprule
时刻 & 状态 & 关键边界 \\
\midrule
exec 提交 & 新 mm 已安装，代码段只有 VMA 或部分 PTE & 旧程序不能继续运行 \\
第一次取指 & 入口页 PTE 不 present，CPU 产生缺页异常 & faulting 指令还没执行 \\
缺页处理 & 根据 VMA 找到 ELF 文件和页偏移 & 必须检查 VMA 可执行权限 \\
读入页面 & page cache miss 时等待存储 I/O & 当前 task 可睡眠，调度器运行别人 \\
安装 PTE & PTE 标成 user、present、read、exec、not writable & 权限不能超过 ELF 段权限 \\
重试取指 & 同一 RIP 再次执行，这次翻译成功 & demand paging 对用户程序透明 \\
\bottomrule
\end{longtable}

这条路径解释了两个现象。第一，大程序 exec 可以很快返回到用户态，因为只物化了马上需要的页；后面真正访问到的代码和数据才陆续缺页。第二，exec 已经成功不代表后续每次取指和读数据都不会失败：底层文件系统 I/O 错误、映射被异常截断、权限标志错误，都可能在第一次访问某页时暴露。教材若只说“exec 装载 ELF”，读者会误以为装载一定是一次性读完整个文件；真实边界是“建立地址空间契约”和“按需把页物化”两步。

\subsection{用户栈构造}

用户栈是 ABI 契约的一部分。教学系统可以简化 auxv，但不能把 argv 字符串和 argv 指针混为一谈。

\begin{lstlisting}
build_user_stack(argv, envp, auxv):
  sp = USER_STACK_TOP
  for s in reverse(strings(argv, envp)):
    sp -= len(s) + 1
    copy_to_new_user(sp, s)
    remember_pointer(s, sp)
  sp = align_down(sp, 16)
  push_null()
  push_auxv(auxv)
  push_null()
  push_pointers(envp)
  push_null()
  push_pointers(argv)
  push(argc)
  return sp
\end{lstlisting}

栈构造要设置总大小上限，防止用户传入巨量参数耗尽内核内存。所有用户字符串必须先复制到内核或临时页中，因为在 \code{exec} 期间旧地址空间即将被替换；若边复制边销毁旧 mm，就可能读到无效地址。

\subsection{提交点}

提交阶段要尽量短，并且顺序清晰：

\begin{lstlisting}
commit_exec(new_image):
  old_mm = current.mm
  current.mm = new_image.mm
  current.regs.ip = new_image.entry
  current.regs.sp = new_image.stack_pointer
  reset_caught_signal_handlers(current)
  close_cloexec_fds(current.files)
  activate_mm(new_image.mm)
  drop_mm(old_mm)
  return_to_user()
\end{lstlisting}

提交点之后不应再执行可能失败且需要返回旧程序的操作。比如关闭 \code{CLOEXEC} fd 一般可以在提交逻辑中完成，因为关闭失败不应阻止 exec 成功；但读取解释器、分配栈页和权限检查必须放在提交点之前。

\section{提交点、失败回滚与身份变化}

\begin{enumerate}
\tightlist
\item 正常 ELF：静态程序能打印 argc/argv/envp，入口地址和栈对齐正确。
\item 坏 magic：返回 \code{ENOEXEC}，旧程序继续运行。
\item 无执行权限：返回 \code{EACCES}，fd table 不变化。
\item 段溢出：构造 \code{p\_vaddr + p\_memsz} 溢出的 ELF，装载器必须拒绝。
\item BSS 清零：全局未初始化数组初值必须为零。
\item \code{FD\_CLOEXEC}：带标志 fd 在新程序中不可见，普通 fd 仍可用。
\item 大 argv：超过限制返回 \code{E2BIG}，不破坏旧地址空间。
\item 动态链接入口：带 \code{PT\_INTERP} 时入口应进入解释器而不是主程序 \code{e\_entry}。
\end{enumerate}

\begin{rubricbox}
掌握本章内容后，应当能够区分 fork 与 exec，解释 ELF program header 的运行时意义，
画出初始用户栈布局，并指出 exec 的失败回滚边界和提交点。
\end{rubricbox}

\subsection{映像替换涉及的完整状态集合}

\code{exec} 很容易被误讲成“把一个 ELF 读进内存”。真正的主线更宽：它同时替换地址空间、
重建用户栈、改变入口寄存器、裁剪 fd 继承、可能改变凭据、处理多线程、交给解释器或动态
链接器，并把旧程序的错误恢复边界切断。把这些细节归入同一张状态账本，才能避免只把装载
理解为几次内存映射。

\begin{longtable}{p{0.20\linewidth}p{0.38\linewidth}p{0.30\linewidth}}
\toprule
账本 & exec 要修改什么 & 不能出错的边界 \\
\midrule
地址空间 & 新 \code{mm}、VMA、栈、brk、mmap base、VDSO & 提交前失败必须保留旧 \code{mm}，提交后不能再回旧代码 \\
文件与路径 & 可执行文件、解释器、脚本原文件、mount flags & 权限检查必须绑定稳定 inode/file，而不是路径字符串 \\
fd 继承 & 普通 fd 继承，\code{FD\_CLOEXEC} fd 关闭 & 不能在 exec 失败时提前关闭旧程序 fd \\
凭据与安全 & setuid、file capability、LSM、no\_new\_privs、dumpable & 提权 exec 要清理危险环境和调试关系 \\
线程组 & 成功后只保留调用 exec 的线程 & 旧用户态锁和其他线程栈不能进入新映像 \\
用户态启动 & \code{argc/argv/envp/auxv}、入口 RIP、栈对齐 & 动态链接器和 libc 在 \code{main} 前就依赖这些字段 \\
\bottomrule
\end{longtable}

这张账本能解释很多看似零散的规则。脚本 \code{\#!/bin/sh} 需要重新构造 argv，是因为真正执行的是解释器；动态 ELF 入口先到 \code{ld-linux}，是因为重定位和共享库加载属于用户态运行时；\code{no\_new\_privs} 会压制提权，是因为 seccomp 和沙箱不能允许“先收紧入口、再 exec 一个 setuid 程序拿回更大权限”；多线程 exec 要清掉其他线程，是因为旧程序里的用户态 mutex、TLS、栈和信号现场对新程序没有可解释语义。

再用一个带提权和动态链接的例子收束。进程执行一个 setuid root 的动态链接程序时，内核不能只检查 ELF magic。它要先解析可执行 inode，计算新凭据，判断 \code{no\_new\_privs}、文件 capability 和 LSM 策略；若发生权限边界变化，运行时还要进入 secure mode，忽略或清理会影响动态链接器的环境变量，例如 \code{LD\_PRELOAD}。否则攻击者可以让高权限程序在启动前加载自己控制的共享库。

\begin{lstlisting}
privileged_exec_boundary(file):
  old_cred = current.cred
  new_cred = compute_exec_cred(file, old_cred)
  if privilege_will_change(old_cred, new_cred):
    enable_secureexec()
    suppress_unsafe_loader_environment()
    restrict_dump_and_ptrace()
  install_new_image_and_cred_at_commit()
\end{lstlisting}

这里的提交点必须把“新程序身份”和“新程序映像”作为一组状态安装。若凭据已经提升但地址空间仍是旧程序，旧代码就拿到了不该有的身份；若地址空间已经换成新程序但 fd 和环境还没裁剪，新程序又可能短暂看到不该继承的能力。好的教学实现即使简化 setuid 和动态链接，也要保留这条不变量：exec 成功后看到的是一组自洽的新运行身份，exec 失败后旧程序仍保持原身份和原资源。

\begin{rubricbox}
分析 \code{exec} 时，至少要区分四个点：仍可回滚到旧程序的点、开始销毁旧用户映像的点、
新凭据和新地址空间一起生效的点、返回用户态且永不返回旧代码的点。缺少这四个点，就很难
判断错误路径是否安全。
\end{rubricbox}

\section{动态链接、解释器脚本与辅助向量}

\subsection{ELF 结构字段}

分析 ELF 装载时要区分三类头：\code{Elf64\_Ehdr} 描述整个文件，\code{Elf64\_Phdr}
描述运行时段，\code{Elf64\_Shdr} 描述链接或调试所用的节。内核运行装载主要依据
program header；section header 即使被移除，也不影响程序运行。

\begin{longtable}{p{0.22\linewidth}p{0.34\linewidth}p{0.32\linewidth}}
\toprule
字段或类型 & 含义 & 装载用途 \\
\midrule
\code{PT\_LOAD} & 可映射段 & 建立 VMA，设置 R/W/X 权限 \\
\code{PT\_INTERP} & 动态解释器路径 & 装载 \code{ld-linux.so} 并把入口交给它 \\
\code{PT\_PHDR} & program header 自身位置 & 传给用户态动态链接器 \\
\code{PT\_GNU\_STACK} & 栈权限提示 & 决定用户栈是否可执行 \\
\code{e\_entry} & ELF 入口虚拟地址 & 静态程序入口或解释器最终跳转参考 \\
\bottomrule
\end{longtable}

Linux 的 ELF 装载路径同样遵循清晰阶段：读取文件头和程序头，检查解释器，准备新地址空间，
映射可装载段，建立堆边界，构造用户栈和辅助向量，最后设置返回用户态所需的寄存器现场。

\subsection{aux vector 布局}

auxv 是内核放到用户栈上的键值数组，给动态链接器和 libc 传递信息。常见项包括 \code{AT\_PHDR}、\code{AT\_PHENT}、\code{AT\_PHNUM}、\code{AT\_ENTRY}、\code{AT\_BASE}、\code{AT\_PAGESZ}、\code{AT\_RANDOM}、\code{AT\_EXECFN}、\code{AT\_NULL}。

\begin{lstlisting}
user stack top:
  strings for argv/envp
  random bytes
  auxv: AT_PHDR, value
        AT_PHENT, value
        AT_PHNUM, value
        AT_ENTRY, value
        AT_RANDOM, pointer
        AT_NULL, 0
  envp pointers, NULL
  argv pointers, NULL
  argc
\end{lstlisting}

\code{AT\_RANDOM} 为 stack canary 和 ASLR 提供随机种子；\code{AT\_BASE} 指向动态链接器基址；\code{AT\_PHDR} 让动态链接器找到主程序段表。若 auxv 构造错误，程序可能在 \code{main} 前就崩溃。

\subsection{解释器脚本和 binfmt}

\code{\#! /bin/sh} 脚本不是 ELF。内核识别 shebang 后，把解释器路径和脚本路径重写成新的 argv，再递归执行解释器。Linux 用 binfmt 框架支持多种二进制格式：ELF、脚本、misc 注册格式等。

\begin{lstlisting}
execve("script.sh", ["script.sh", "a"]):
  read first line "#!/bin/sh"
  new argv = ["/bin/sh", "script.sh", "a"]
  execve("/bin/sh", new argv)
\end{lstlisting}

这里要限制递归深度，否则脚本解释器互相指向会无限递归。binfmt\_misc 允许用户注册自定义格式，例如通过魔数把某类文件交给解释器。安全检查必须覆盖最终解释器和原脚本，不能只检查第一层路径。

\subsection{vfork + exec 为什么高效但危险}

\code{vfork} 让子进程在 exec 或 exit 前与父进程共享地址空间，父进程通常被挂起。它避免了复制页表和 COW 设置，适合立刻 exec 的路径；但子进程若在 exec 前修改内存或调用复杂函数，会破坏父进程状态。

\begin{lstlisting}
child = vfork()
if child == 0:
  // only async-signal-safe, minimal setup
  dup2(fd, STDOUT_FILENO)
  execve(path, argv, envp)
  _exit(127)
\end{lstlisting}

现代系统常用 \code{posix\_spawn} 在库和内核之间选择更安全高效的实现。理解 \code{vfork} 的意义，不是鼓励随便使用，而是理解 exec 前后地址空间提交点的重要性。

\subsection{setuid、凭据变化与 no\_new\_privs}

exec 不只是换代码。执行 setuid/setgid 文件时，进程凭据可能变化；capability 也会根据文件能力、bounding set、ambient set 和 no\_new\_privs 重新计算。动态链接器和环境变量处理必须考虑提权执行，避免用户用 \code{LD\_PRELOAD} 一类环境影响高权限程序。

\begin{lstlisting}
exec_security_transition(file):
  read inode mode, owner, file capabilities
  if no_new_privs:
    suppress privilege gain
  compute new effective/permitted/inheritable caps
  clear dangerous personality and env effects
  install new credentials at commit
\end{lstlisting}

这解释了为什么 exec 路径会调用 LSM hook 和凭据提交逻辑。安全检查不能只看 ELF 格式合法，还要看执行后主体身份是否改变、哪些 fd 继承、dumpable 标志如何更新、ptrace 是否允许继续。

\subsection{ASLR、PIE 与栈保护输入}

exec 建立新地址空间时，会选择随机化的 mmap base、stack base、brk 和 VDSO 位置。PIE 主程序也可随机基址加载。内核还向 auxv 提供 \code{AT\_RANDOM}，供 libc 初始化 stack canary。

\begin{longtable}{p{0.24\linewidth}p{0.48\linewidth}}
\toprule
机制 & exec 阶段输入 \\
\midrule
ASLR & 随机选择 mmap、stack、brk、VDSO、PIE 基址 \\
stack canary & \code{AT\_RANDOM} 提供随机字节 \\
NX stack & \code{PT\_GNU\_STACK} 和默认策略决定栈可执行性 \\
RELRO/动态链接 & program header 和 dynamic section 给链接器信息 \\
\bottomrule
\end{longtable}

这些保护不是 exec 独立完成的，但 exec 负责把初始地址布局和 auxv 种子交给用户态运行时。若地址布局固定或随机数可预测，后续保护会明显削弱。

\subsection{exec 与多线程}

多线程进程调用 exec 时，成功后只留下调用 exec 的线程，其他线程被清理。因为新程序映像不可能继承旧程序的线程栈、锁状态和用户态同步对象。内核必须停止同一线程组中的其他线程，并让共享资源进入一致状态。

\begin{lstlisting}
de_thread_for_exec(current):
  signal other threads to exit
  wait until thread group is single-threaded
  take over thread group leader identity if needed
  proceed to install new mm
\end{lstlisting}

这也是为什么 exec 提交点必须非常谨慎：旧线程可能持有用户态锁，但 exec 后这些锁没有意义；内核对象如 fd table、credentials、signal struct 还要按规则继承或重置。

\subsection{binfmt\_misc 与解释器安全边界}

binfmt 框架允许注册自定义格式，把某些 magic 或扩展名交给解释器。它方便运行脚本、模拟器格式或语言运行时，但也扩大了 exec 策略面。解释器路径、参数重写、递归深度、权限继承和 mount noexec 都必须纳入检查。

\begin{lstlisting}
registered format:
  magic = "\xCA\xFE"
  interpreter = "/usr/bin/custom-runner"
exec file:
  kernel rewrites argv:
    custom-runner original-file original-args...
\end{lstlisting}

安全模型必须说明权限检查作用于原文件、解释器，还是二者都检查。只检查脚本可执行而不检查解释器，或绕过 noexec 挂载，都会形成边界漏洞。
